% !TEX program = xelatex
\documentclass[11pt,a4paper]{article}
\usepackage[a4paper, left=2.5cm, right=4cm, top=2.5cm, bottom=2cm]{geometry}
\usepackage{setspace}
\onehalfspacing
\usepackage{fontspec}
\setmainfont{Times New Roman}
\usepackage{polyglossia}
\setmainlanguage[spelling=new]{german}
\usepackage{amsmath,amssymb,mathtools}
\usepackage{booktabs}
\usepackage{tabularx}
\usepackage{array}
\usepackage{graphicx}
\usepackage{float}
\usepackage{enumitem}
\usepackage{hyperref}
\usepackage{xcolor}
\usepackage{caption}
\usepackage{parskip}
\renewcommand{\footnotesize}{\fontsize{10pt}{12pt}\selectfont}

\captionsetup[table]{
    position=above,
    justification=raggedright,
    singlelinecheck=false,
    labelfont=bf,
    labelsep=newline,
    textfont=it
}
\captionsetup[figure]{
    position=below,
    justification=raggedright,
    singlelinecheck=false,
    labelfont=bf,
    labelsep=newline,
    textfont=it
}

\hypersetup{
  colorlinks=true,
  linkcolor=blue!70!black,
  citecolor=blue!70!black,
  urlcolor=blue!70!black,
  pdftitle={Stochastische ROI-Bewertung mittels Monte-Carlo-Simulation},
  pdfauthor={Elias Scho \& Fabian Finkenzeller},
}

\title{\textbf{Monte-Carlo-Simulation zur stochastischen Bewertung von Investitionsrenditen in der \"Olindustrie}}
\author{Elias Scho \& Fabian Finkenzeller}
\date{Juni 2026}
\begin{document}

% ==================================================================
% TITELSEITE
% ==================================================================
\maketitle
\thispagestyle{empty}
\begin{center}
\begin{tabular}{ll}
\toprule
\textbf{Feld}    & \textbf{Wert} \\
\midrule
Titel       & Stochastische ROI-Bewertung mittels Monte-Carlo-Simulation \\
Typ         & Wissenschaftliches Working Paper \\
Version     & 3.0 \\
Datum       & Juni 2026 \\
\bottomrule
\end{tabular}
\end{center}

\vspace{1.5cm}

\begin{abstract}\noindent
Die \"Olindustrie zeichnet sich durch extreme Kapitalintensit\"at und hohe Unsicherheit aus. Herk\"ommliche deterministische Investitionsrechnungen, die auf Einzelpunktsch\"atzungen basieren, bilden die realen Projektrisiken systematisch unzureichend ab. Diese Arbeit entwickelt ein stochastisches ROI-Modell, das die Unsicherheiten \"olindustrieller Investitionsentscheidungen methodisch quantifiziert. Das Modell simuliert vier Kernvariablen mittels Monte-Carlo-Simulation: CAPEX, OPEX, F\"ordervolumen und \"Olpreis. Die resultierende ROI-Verteilung liefert nicht nur einen Erwartungswert, sondern ein vollst\"andiges Risikoprofil inklusive Value at Risk und Verlustwahrscheinlichkeit.
\end{abstract}

\newpage
\tableofcontents
\newpage

% ==================================================================
% ABSCHNITT 1: EINLEITUNG
% ==================================================================
\section{Einleitung}

\subsection{Motivation}

Die \"Olindustrie zeichnet sich durch eine einzigartige Kombination von extremer Kapitalintensit\"at und hoher Unsicherheit aus, begleitet von langen Investitionszeitr\"aumen. Einzelne Projekte wie Offshore-Bohrungen, Frackingvorhaben oder Pipelineinfrastruktur erfordern Investitionen im Bereich von Hunderten Millionen bis Milliarden Euro, bevor auch nur ein einziges Barrel gef\"ordert wird. Gleichzeitig unterliegen die treibenden Erfolgsfaktoren, n\"amlich Rohstoffpreise, geologische Reserven, Betriebskosten und regulatorische Rahmenbedingungen, einer Volatilit\"at, die herk\"ommliche Planungsmethoden systematisch \"uberfordert.

Die Finanzkrise 2008, der \"Olpreis-Crash 2014--2016 und die COVID-bedingte Preiskatastrophe 2020 haben eindrucksvoll gezeigt, dass deterministische Planungsans\"atze, die auf Einzelpunktsch\"atzungen (\textit{Point Estimates}) basieren, die Realit\"at systematisch verfehlen. Ein Projekt, das bei einem \"Olpreis von \$80/Barrel hervorragend rentabel erscheint, kann bei \$40/Barrel katastrophal sein.

\subsection{Problemstellung}

Herk\"ommliche deterministische Investitionsrechnungen basieren typischerweise auf punktuellen Sch\"atzwerten f\"ur \"Olpreis, F\"ordermengen und Kosten. Sie verdichten das Ergebnis zu einem einzigen ROI-Wert, der eine irref\"uhrende Gewissheit \"uber die Projektrentabilit\"at suggeriert. Dieser Ansatz ignoriert drei kritische Aspekte:

\begin{enumerate}[label=\arabic*., itemsep=4pt]
  \item \textbf{Unsicherheit der Inputgr\"o{\ss}en}: Keine der vier Kernvariablen (CAPEX, OPEX, F\"ordervolumen, \"Olpreis) ist mit Sicherheit vorhersehbar.
  \item \textbf{Asymmetrische Risikoprofile}: Die Kernvariablen folgen nicht einheitlich einer symmetrischen Normalverteilung, sondern weisen individuelle, oft asymmetrische Risikoprofile auf.
  \item \textbf{Nicht-lineare Aggregation}: Die Wechselwirkungen zwischen den Variablen werden in deterministischen Modellen nicht abgebildet.
\end{enumerate}

Besonders relevant ist diese Problematik f\"ur investitionsbasierte Antragsdaten (Ex-ante-Daten). Da Investitionsentscheidungen \"ublicherweise auf reinen Prognosen und Sch\"atzwerten beruhen, kommt es bei statischen Modellen fast unweigerlich zu systematischen Verzerrungen bei der Risikobewertung.

\subsection{Zielsetzung}

Das Ziel dieses Papers ist die Entwicklung und Darstellung eines stochastischen ROI-Modells, das die Unsicherheiten \"olindustrieller Investitionsentscheidungen methodisch sauber quantifiziert. Konkret werden wir:

\begin{itemize}[itemsep=2pt]
  \item Ein mathematisch fundiertes Vier-Variablen-Modell entwickeln, das CAPEX, OPEX, F\"ordervolumen und \"Olpreis als Zufallsvariablen modelliert
  \item Die Wahl der Wahrscheinlichkeitsverteilungen wissenschaftlich begr\"unden
  \item Eine Monte-Carlo-Simulation implementieren, die die aggregierte ROI-Verteilung berechnet
  \item Die Limitationen des Ansatzes kritisch reflektieren
\end{itemize}

\newpage

% ==================================================================
% ABSCHNITT 2: THEORETISCHER HINTERGRUND
% ==================================================================
\section{Theoretischer Hintergrund}

\subsection{Investitionsrechnung unter Unsicherheit}

Die klassische Investitionsrechnung kennt mehrere Verfahren zur Bewertung von Investitionen: Statische Verfahren (Rentabilit\"atsrechnung, Amortisationsrechnung) und dynamische Verfahren (Kapitalwertmethode, Interner Zinsfu{\ss}, Annuit\"atenmethode). Allen gemeinsam ist, dass sie in der Regel mit deterministischen Cashflows arbeiten.

Der Return on Investment (ROI) berechnet sich gem\"a{\ss}:
\begin{equation}
\text{ROI} = \frac{\text{Gewinn}}{\text{Investition}} \times 100\%
\end{equation}

bzw.\ in der \"olindustriellen Spezifikation:
\begin{equation}
\text{ROI} = \frac{\text{{\"O}lpreis} \times \text{F\"ordervolumen} - \text{CAPEX} - \text{OPEX}_{\text{ges}}}{\text{CAPEX}}
\end{equation}

wobei $\text{OPEX}_{\text{ges}} = \text{OPEX}_{\text{j\"ahrlich}} \times T$ das Produkt aus j\"ahrlichen Betriebskosten und der Projektlaufzeit $T$ in Jahren ist.

Durch die Modellierung der Eingabevariablen als Zufallsvariablen wird der ROI selbst zu einer Zufallsvariable. Der ROI ist dadurch keine Zahl mehr, sondern eine Verteilung -- und die Form dieser Verteilung enth\"alt die entscheidenden Informationen \"uber das Risiko.

\subsection{Ex-ante vs.\ Ex-post Steuerung}

Ein zentraler methodischer Punkt unseres Ansatzes betrifft die Unterscheidung zwischen Ex-ante- und Ex-post-Perspektive:

\textbf{Ex-ante (vorab)}: Die Investitionsentscheidung wird auf Basis von Antragsdaten getroffen -- Sch\"atzungen von Ingenieuren, geologischen Gutachten und Marktanalysen. Zum Zeitpunkt der Entscheidung existieren keine tats\"achlichen Ertr\"age, nur Prognosen.

\textbf{Ex-post (nachtr\"aglich)}: Nach Projektabschluss k\"onnen die tats\"achlichen Kosten und Ertr\"age gemessen werden.

Unser Modell ist rein \textbf{Ex-ante} konzipiert: Es modelliert die Investitionsentscheidung aus der Perspektive des Bewertenden zum Zeitpunkt $t_0$, bevor das Projekt startet. Dies bedeutet insbesondere, dass wir von einer \glqq Fire-and-Forget\grqq{}-Investition ausgehen: Nach Bewilligung des Antrags erfolgen keine weiteren adaptiven Eingriffe (Realoptionen werden nicht modelliert -- siehe Limitationen in Abschnitt~\ref{sec:limitationen}).

\subsection{Funktionsweise der Monte-Carlo-Simulation}

Die Monte-Carlo-Simulation (MCS) ist eine rechnergest\"utzte Methode zur numerischen L\"osung von Problemen unter Unsicherheit. Ihre Urspr\"unge liegen in den Arbeiten von \textbf{Stanislaw Ulam} und \textbf{John von Neumann} am Los Alamos National Laboratory in den 1940er Jahren, wo sie zur Modellierung von Neutronendiffusion entwickelt wurde.

Den \"Ubergang in die betriebswirtschaftliche Investitionsrechnung vollzog \textbf{David B.~Hertz} mit seinem wegweisenden Artikel \textit{``Risk Analysis in Capital Investment''} \cite{hertz1964}. Hertz demonstrierte erstmals, wie probabilistische Simulation die traditionelle Einzelpunktsch\"atzung bei Investitionsentscheidungen abl\"osen kann.

In der Erd\"olindustrie wurde die Methode ma{\ss}geblich durch \textbf{Paul D.~Newendorp} etabliert, dessen Standardwerk \textit{``Decision Analysis for Petroleum Exploration''} \cite{newendorp2000} die probabilistische Entscheidungsanalyse f\"ur Exploration und Produktion (E\&P) systematisierte und die Dreiecksverteilung als Standardwerkzeug f\"ur \textit{Expert Elicitation} begr\"undete.

Die MCS basiert auf dem Gesetz der gro{\ss}en Zahlen und folgt einem vierstufigen Prozess:

\begin{enumerate}[label=\arabic*., itemsep=4pt]
  \item \textbf{Definition stochastischer Inputvariablen}: Jede der vier Variablen (CAPEX, OPEX, F\"ordervolumen, \"Olpreis) wird als Zufallsvariable mit einer spezifischen Wahrscheinlichkeitsverteilung modelliert.
  \item \textbf{Zuf\"alliges Ziehen (Sampling)}: F\"ur jede der $N$ Iterationen werden f\"ur alle vier Variablen zuf\"allige Werte aus ihren jeweiligen Verteilungen gezogen.
  \item \textbf{Berechnung des ROI}: F\"ur jeden Satz gezogener Werte wird der ROI nach der definierten Formel berechnet.
  \item \textbf{Aggregation}: Nach $N$ Iterationen ergibt sich eine Verteilung von $N$ ROI-Werten, die als Histogramm dargestellt werden kann.
\end{enumerate}

Der zentrale Erkenntnisgewinn liegt darin, dass die ROI-Verteilung nicht nur einen Erwartungswert liefert, sondern die gesamte Risikobandbreite zeigt -- einschlie{\ss}lich der Wahrscheinlichkeit, dass das Projekt Verluste macht ($\text{ROI} < 0$).

% ==================================================================
% ABSCHNITT 3: METHODIK UND MODELLDESIGN
% ==================================================================
\section{Methodik und Modelldesign}

\subsection{Variablenselektion}

Die Vier-Variablen-Struktur unseres Modells wurde bewusst gew\"ahlt. In der \"olindustriellen Investitionsbewertung dominieren diese vier Gr\"o{\ss}en das Risikoprofil eines Projekts:

\paragraph{1. CAPEX (Capital Expenditure -- Investitionskosten)}
\begin{itemize}[itemsep=1pt]
  \item Umfasst: Bohrkosten, Plattformbau, Pipeline-Infrastruktur, Genehmigungskosten
  \item Charakteristikum: Einmalige, hohe Investition zu Projektbeginn
  \item \"Uberraschungen: Kosten\"uberschreitungen von 20--50\% sind branchen\"ublich \cite{spe35990}
\end{itemize}

\paragraph{2. OPEX (Operational Expenditure -- Betriebskosten)}
\begin{itemize}[itemsep=1pt]
  \item Umfasst: Wartung, Personal, Logistik, Umweltschutzma{\ss}nahmen
  \item Charakteristikum: Laufende Kosten \"uber die gesamte F\"orderdauer
  \item \"Uberraschungen: Kann durch technische Probleme oder regulatorische \"Anderungen stark variieren
\end{itemize}

\paragraph{3. F\"ordervolumen (Reserven / Produktion)}
\begin{itemize}[itemsep=1pt]
  \item Umfasst: Die tats\"achlich f\"orderbare \"Olmenge \"uber die Projektlebensdauer
  \item Charakteristikum: Geologische Unsicherheit -- selbst nach Explorationsbohrungen bleiben signifikante Unsicherheiten
  \item \"Uberraschungen: Reserven k\"onnen um 30--50\% von der Sch\"atzung abweichen
\end{itemize}

\paragraph{4. \"Olpreis (Marktpreis)}
\begin{itemize}[itemsep=1pt]
  \item Umfasst: Der durchschnittliche Verkaufspreis pro Barrel \"uber die Projektdauer
  \item Charakteristikum: Exogene Marktvariable, nicht kontrollierbar
  \item \"Uberraschungen: Historische Volatilit\"at von $\pm$40\% \cite{leonteq2023} innerhalb weniger Jahre
\end{itemize}

Weitere Faktoren (Steuern, Diskontierungss\"atze, Devisenkurse) spielen ebenfalls eine Rolle, aber die vier genannten Variablen erkl\"aren den Gro{\ss}teil der ROI-Varianz und halten das Modell interpretierbar.

\subsection{Verteilungsannahmen}

Die Wahl der Wahrscheinlichkeitsverteilungen ist der wissenschaftlich kritischste Teil des Modells. Jede Verteilung muss zwei Kriterien erf\"ullen: (1)~Sie muss die Natur der Unsicherheit sachgerecht abbilden, und (2)~sie muss mit den verf\"ugbaren Daten (Antragsdaten, Expertensch\"atzungen) kalibrierbar sein.

\subsubsection{Dreiecksverteilung (CAPEX, OPEX, F\"ordervolumen)}

Eine Zufallsvariable $X$ folgt einer Dreiecksverteilung $\text{Tri}(a, b, c)$, wenn ihre Dichtefunktion durch ein Dreieck mit den Eckpunkten $a$ (Minimum), $b$ (Modus/wahrscheinlichster Wert) und $c$ (Maximum) gegeben ist.

Dichte:
\begin{equation}
f(x) = \begin{cases}
  \dfrac{2(x-a)}{(b-a)(c-a)} & \text{f\"ur } a \leq x \leq b \\[8pt]
  \dfrac{2(c-x)}{(c-b)(c-a)} & \text{f\"ur } b < x \leq c \\[4pt]
  0 & \text{sonst}
\end{cases}
\end{equation}

Erwartungswert:
\begin{equation}
E[X] = \frac{a + b + c}{3}
\end{equation}

Varianz:
\begin{equation}
\text{Var}[X] = \frac{a^2 + b^2 + c^2 - ab - ac - bc}{18}
\end{equation}

\textbf{Warum Dreiecksverteilung?}

\begin{enumerate}[label=\arabic*., itemsep=4pt]
  \item \textbf{Datenrealit\"at}: In der Praxis werden Investitionsantr\"age selten mit vollst\"andigen historischen Datens\"atzen versehen. Ingenieure und Geologen liefern typischerweise drei Werte: einen Worst-Case (Minimum), einen Best-Case (Maximum) und einen Most-Likely-Case (Modus). Genau diese drei Parameter definieren die Dreiecksverteilung \cite{newendorp2000}.
  \item \textbf{Wissenschaftliche Begr\"undung}: Die Dreiecksverteilung erzwingt keine Symmetrie und hat einen endlichen Support, was negative Werte bei Kosten ausschlie{\ss}t. Sie ist die minimale Informationsverteilung, die mit den drei gegebenen Sch\"atzwerten konsistent ist.
  \item \textbf{Industriestandard}: Die Society of Petroleum Engineers (SPE) empfiehlt in SPE~35990 \cite{spe35990} explizit die Dreiecksverteilung f\"ur die probabilistische Sch\"atzung von Bohrkosten und CAPEX-Komponenten, wenn keine ausreichende historische Datenbasis f\"ur eine parametrische Verteilungsanpassung vorliegt.
\end{enumerate}

\textbf{Praktische Kalibrierung:}
\begin{itemize}[itemsep=1pt]
  \item CAPEX: $\text{Tri}(500\text{M},\; 750\text{M},\; 1{,}2\text{B})$ -- typische Bandbreite f\"ur Offshore-Projekte
  \item OPEX (j\"ahrlich): $\text{Tri}(80\text{M},\; 120\text{M},\; 200\text{M})$ -- laufende Kosten pro Jahr
  \item F\"ordervolumen: $\text{Tri}(50\text{M},\; 150\text{M},\; 300\text{M})$ -- Barrel \"uber Projektlebensdauer
  \item Projektlaufzeit: $T = 10$ Jahre
  \item OPEX gesamt: $\text{OPEX}_{\text{ges}} = \text{OPEX}_{\text{j\"ahrlich}} \times T$
\end{itemize}

\subsubsection{Lognormalverteilung (\"Olpreis)}

Eine Zufallsvariable $X$ folgt einer Lognormalverteilung $\text{LogN}(\mu, \sigma^2)$, wenn $\ln(X) \sim \mathcal{N}(\mu, \sigma^2)$.

Dichte:
\begin{equation}
f(x) = \frac{1}{x \cdot \sigma \cdot \sqrt{2\pi}} \cdot \exp\!\left(-\frac{(\ln(x) - \mu)^2}{2\sigma^2}\right), \quad x > 0
\end{equation}

Erwartungswert:
\begin{equation}
E[X] = \exp\!\left(\mu + \frac{\sigma^2}{2}\right)
\end{equation}

Varianz:
\begin{equation}
\text{Var}[X] = \bigl(\exp(\sigma^2) - 1\bigr) \cdot \exp(2\mu + \sigma^2)
\end{equation}

\textbf{Warum Lognormalverteilung f\"ur den \"Olpreis?}

\begin{enumerate}[label=\arabic*., itemsep=4pt]
  \item \textbf{Nicht-Negativit\"at}: Die Lognormalverteilung hat nat\"urlich $x > 0$.
  \item \textbf{Rechtsschiefe}: Historische \"Olpreise zeigen eine deutliche Rechtsschiefe -- sie k\"onnen theoretisch beliebig hoch ansteigen, fallen aber nicht unter null.
  \item \textbf{Finanzmathematischer Standard}: Die Lognormalverteilung ist der etablierte Standard in der Finanzmathematik zur Modellierung von Rohstoffpreisen. Brennan und Schwartz \cite{brennan1985} modellierten erstmals Rohstoffinvestitionen mittels geometrischer Brownscher Bewegung, bei der der Preis zu jedem festen Zeitpunkt $T$ lognormalverteilt ist.
  \item \textbf{Multiplikatives Wachstum}: \"Olpreis\"anderungen sind proportional. Proportionale \"Anderungen \"uber die Zeit f\"uhren mathematisch zu Lognormalverteilungen.
\end{enumerate}

\textbf{Kalibrierung aus Marktdaten:}
\begin{itemize}[itemsep=1pt]
  \item Historischer mittlerer \"Olpreis: ${\sim}\$70$/Barrel
  \item Historische Volatilit\"at: $\sigma \approx 0{,}35$
  \item Parameter: $\mu = \ln(70) - \sigma^2/2$, \; $\sigma = 0{,}35$
\end{itemize}

\subsection{Mathematisches Modell}

Das ROI-Modell berechnet sich f\"ur jeden Simulationsschritt $i$ ($i = 1, \dots, N$) wie folgt:

\textbf{Gegeben:}
\begin{align}
  \text{CAPEX}_i         &\sim \text{Tri}(a_c,\, b_c,\, c_c)   & &\text{-- Investitionskosten} \\
  \text{OPEX}_{\text{annual},i}  &\sim \text{Tri}(a_o,\, b_o,\, c_o) & &\text{-- j\"ahrliche Betriebskosten} \\
  \text{OPEX}_{\text{ges},i}     &= \text{OPEX}_{\text{annual},i} \times T & &\text{-- Betriebskosten gesamt} \\
  V_i                    &\sim \text{Tri}(a_v,\, b_v,\, c_v)   & &\text{-- F\"ordervolumen (Barrel)} \\
  P_i                    &\sim \text{LogN}(\mu_p,\, \sigma_p^2) & &\text{-- \"Olpreis (\$/Barrel)}
\end{align}

\textbf{Berechnung:}
\begin{align}
  \text{Umsatz}_i       &= P_i \times V_i \\
  \text{Gesamtkosten}_i &= \text{CAPEX}_i + \text{OPEX}_{\text{ges},i} \\
  \text{Gewinn}_i       &= \text{Umsatz}_i - \text{Gesamtkosten}_i \\
  \text{ROI}_i          &= \frac{\text{Gewinn}_i}{\text{CAPEX}_i}
\end{align}

Nach $N$ Iterationen ergibt sich die empirische ROI-Verteilung $\{\text{ROI}_1, \text{ROI}_2, \dots, \text{ROI}_N\}$, die ausgewertet wird durch:

\begin{itemize}[itemsep=2pt]
  \item \textbf{Erwartungswert (Mean ROI)}: $\displaystyle \bar{R} = \frac{1}{N}\sum_{i=1}^{N} \text{ROI}_i$
  \item \textbf{Standardabweichung (Volatilit\"at)}: $\displaystyle \hat{\sigma} = \sqrt{\frac{1}{N-1}\sum_{i=1}^{N}(\text{ROI}_i - \bar{R})^2}$
  \item \textbf{Value at Risk (VaR)}: $P(\text{ROI} \leq \text{VaR}_\alpha) = \alpha$, typischerweise $\alpha = 5\%$
  \item \textbf{Probability of Loss}: $P(\text{ROI} < 0)$ -- Anteil der Iterationen mit negativem ROI
  \item \textbf{Median (P50)}: Der ROI-Wert, bei dem 50\% der Simulationen dar\"uber und 50\% darunter liegen
\end{itemize}

\subsection{Simulationsparameter}

\begin{table}[H]
\centering
\begin{tabular}{lll}
\toprule
\textbf{Parameter} & \textbf{Standardwert} & \textbf{Begr\"undung} \\
\midrule
Anzahl Iterationen ($N$) & 10.000 & Standardkonvergenz \\
Zufallszahlengenerator & Mersenne Twister & Reproduzierbar (Seed 12345) \\
Korrelationen & Keine (unabh\"angig) & Basis-Modell \\
Projektlaufzeit $T$ & 10 Jahre & Typische Offshore-Projektdauer \\
\bottomrule
\end{tabular}
\caption{Simulationsparameter und ihre Begr\"undung}
\end{table}

\subsection{Hinweis zur Modellkalibrierung}

Die in diesem Paper verwendeten Parameterwerte sind als illustrative Beispielwerte zu verstehen. Bei der konkreten Kalibrierung des Modells zeigte sich, dass die aktuellen Parametersch\"atzungen zu einem extrem hohen erwarteten ROI f\"uhren. Dies liegt daran, dass die j\"ahrlichen OPEX-Sch\"atzungen im Verh\"altnis zum Umsatz (\"Olpreis $\times$ F\"ordervolumen) sehr niedrig ausfallen. In der Praxis sind die Betriebskosten bei Offshore-Projekten deutlich h\"oher, und die hier verwendeten Werte dienen prim\"ar der Demonstration der Methodik. F\"ur eine reale Anwendung m\"ussen die Verteilungsparameter durch umfassende historische Datenanalyse und Expertenvalidierung kalibriert werden.

% ==================================================================
% ABSCHNITT 4: SIMULATION UND ERGEBNISSE
% ==================================================================
\section{Simulation und Ergebnisse}

\subsection{Beschreibung des Datensatzes}

F\"ur die vorliegende Analyse verwenden wir einen synthetischen Datensatz, der auf typischen Antragswerten f\"ur ein mittleres Offshore-{\"O}lf\"orderprojekt basiert.

\paragraph{CAPEX-Parameter (Dreiecksverteilung):}
\begin{itemize}[itemsep=1pt]
  \item Minimum (Worst Case): \$500M
  \item Modus (Most Likely): \$750M
  \item Maximum (Best Case): \$1.200M
\end{itemize}

\paragraph{OPEX-Parameter (Dreiecksverteilung, j\"ahrlich):}
\begin{itemize}[itemsep=1pt]
  \item Minimum: \$80M/Jahr
  \item Modus: \$120M/Jahr
  \item Maximum: \$200M/Jahr
  \item Projektlaufzeit: $T = 10$ Jahre $\rightarrow$ Gesamt-OPEX: $\text{Tri}(800\text{M},\; 1{,}2\text{B},\; 2{,}0\text{B})$
\end{itemize}

\paragraph{F\"ordervolumen-Parameter (Dreiecksverteilung):}
\begin{itemize}[itemsep=1pt]
  \item Minimum: 50M Barrel
  \item Modus: 150M Barrel
  \item Maximum: 300M Barrel
\end{itemize}

\paragraph{{\"O}lpreis-Parameter (Lognormalverteilung):}
\begin{itemize}[itemsep=1pt]
  \item Mittelwert: ${\sim}\$70$/Barrel
  \item Volatilit\"at: $\sigma \approx 0{,}35$
  \item Parameter: $\mu \approx 4{,}20$, $\sigma \approx 0{,}35$
\end{itemize}

\subsection{Durchf\"uhrung der Simulation}

Die Monte-Carlo-Simulation wurde mit $N = 10.000$ Iterationen durchgef\"uhrt (Seed: 12345, Mersenne Twister). In jedem Schritt wurden:
\begin{enumerate}[label=\arabic*., itemsep=2pt]
  \item Vier Zufallszahlen aus den jeweiligen Verteilungen gezogen
  \item Der ROI nach der in Abschnitt~3.3 definierten Formel berechnet
  \item Das Ergebnis der ROI-Stichprobe hinzugef\"ugt
\end{enumerate}

Die Konvergenz wurde durch einen Split-Halftest verifiziert: Die ROI-Verteilung der ersten 5.000 Iterationen weicht von der der zweiten 5.000 Iterationen um weniger als 2\% im Erwartungswert ab.

\subsection{Ergebnisse}

\textbf{Kennzahlen der ROI-Verteilung:}

\begin{table}[H]
\centering
\begin{tabular}{ll}
\toprule
\textbf{Kennzahl} & \textbf{Wert} \\
\midrule
Erwartungswert (Mean ROI)      & $1{.}208\%$ \\
Median (P50)                   & $1{.}052\%$ \\
Standardabweichung              & $759\%$ \\
Value at Risk (5\%)             & $284\%$ \\
Probability of Loss ($P(\text{ROI}<0)$) & $0{,}1\%$ \\
Maximum ROI                    & $6{.}466\%$ \\
Minimum ROI                    & $-140\%$ \\
\bottomrule
\end{tabular}
\caption{Kennzahlen der ROI-Verteilung der Monte-Carlo-Simulation ($N = 10.000$)}
\end{table}

Die extrem hohen ROI-Werte sind auf die Kalibrierung der Eingabeparameter zur\"uckzuf\"uhren: Bei einem mittleren \"Olpreis von \$70/Barrel und einem mittleren F\"ordervolumen von 150M Barrel ergibt sich ein Umsatz von \$10,5 Mrd., dem Gesamtkosten von lediglich ca.\ \$1,95 Mrd.\ (CAPEX \$750M + OPEX \$1,2 Mrd.) gegen\"uberstehen. Die sich daraus ergebende Gewinnmarge ist unrealistisch hoch. Dies unterstreicht die in Abschnitt~3.5 diskutierte Notwendigkeit einer sorgf\"altigen Kalibrierung.

% Platzhalter f\"ur Histogramm-Grafik
\begin{figure}[H]
\centering
\fbox{\parbox{0.85\textwidth}{\centering\vspace{3cm}
\textit{Hinweis: ROI-Histogramm} \\[4pt]
\texttt{figures/roi\_distribution.pdf} -- aus Simulationsdaten generieren
\vspace{3cm}}}
\caption{ROI-Histogramm der Monte-Carlo-Simulation ($N = 10.000$)}
\end{figure}

Die ROI-Verteilung zeigt eine starke Rechtsschiefe: Der Median liegt unter dem arithmetischen Mittel, und die Wahrscheinlichkeit eines Verlustes ist mit $0{,}1\%$ sehr gering.

\subsection{Sensitivit\"atsanalyse}

Um zu verstehen, welche Inputvariable den gr\"o{\ss}ten Einfluss auf die ROI-Streuung hat, wurde eine korrelationsbasierte Sensitivit\"atsanalyse durchgef\"uhrt \cite{spe52949}. F\"ur jede Input-Variable $X_k$ wird der Pearson-Korrelationskoeffizient $r_k$ mit dem ROI berechnet:

\begin{equation}
r(X, Y) = \frac{\sum_{i=1}^{N}(x_i - \bar{x})(y_i - \bar{y})}{\sqrt{\sum_{i=1}^{N}(x_i - \bar{x})^2 \cdot \sum_{i=1}^{N}(y_i - \bar{y})^2}}
\end{equation}

Das Bestimmtheitsma{\ss} $R^2_k = r_k^2$ gibt den Anteil der ROI-Varianz an, der durch die jeweilige Variable erkl\"art wird.

\begin{table}[H]
\centering
\begin{tabular}{llc}
\toprule
\textbf{Variable} & \textbf{Einfluss auf ROI-Varianz} & \textbf{Pearson } $r$ \\
\midrule
{\"O}lpreis        & $51{,}6\%$ & $+0{,}69$ \\
F\"ordervolumen  & $39{,}2\%$ & $+0{,}60$ \\
CAPEX          & $9{,}1\%$  & $-0{,}29$ \\
OPEX           & $0{,}2\%$  & $-0{,}04$ \\
\bottomrule
\end{tabular}
\caption{Sensitivit\"atsanalyse -- Einfluss der Inputvariablen auf die ROI-Varianz}
\end{table}

% Platzhalter f\"ur Tornado-Diagramm
\begin{figure}[H]
\centering
\fbox{\parbox{0.85\textwidth}{\centering\vspace{3cm}
\textit{Hinweis: Tornado-Diagramm} \\[4pt]
\texttt{figures/sensitivity\_tornado.pdf} -- aus Simulationsdaten generieren
\vspace{3cm}}}
\caption{Tornado-Diagramm der Sensitivit\"atsanalyse}
\end{figure}

Die Sensitivit\"atsanalyse zeigt: \"Olpreis und F\"ordervolumen dominieren das Risikoprofil. Der geringe Einfluss der OPEX ist darauf zur\"uckzuf\"uhren, dass die j\"ahrlichen Betriebskosten im Verh\"altnis zum Umsatz sehr niedrig sind. Bei einer realistischeren Kalibrierung w\"are der OPEX-Einfluss deutlich h\"oher.

% ==================================================================
% ABSCHNITT 5: DISKUSSION UND LIMITATIONEN
% ==================================================================
\section{Diskussion und Limitationen}
\label{sec:limitationen}

\subsection{Zusammenhang zwischen den Eingabedaten und der Ergebnisqualit\"at}

Die Qualit\"at der Monte-Carlo-Simulation steht und f\"allt mit der Qualit\"at der Inputdaten. Wenn die Dreiecksverteilungsparameter systematisch verzerrt sind -- z.\,B.\ durch \textit{Optimism Bias} oder \textit{Anchoring} --, dann reflektiert die Output-Verteilung nicht die wahre Unsicherheit, sondern die Verzerrung der Sch\"atzer. Dies wird auch als \glqq Garbage in, Garbage out\grqq{}~bezeichnet.

\textbf{Empfehlung}: Wo m\"oglich, sollten historische Daten zur Validierung der Experten-Sch\"atzungen herangezogen werden. Ein Kalibriertest vergleicht die Vorhersagen vergangener Antr\"age mit den tats\"achlichen Ergebnissen und identifiziert systematische \"Ubersch\"atzungen oder Untersch\"atzungen.

\subsection{Unabh\"angigkeitsannahme}

Unser Basis-Modell geht davon aus, dass die vier Inputvariablen unabh\"angig sind. In der Realit\"at bestehen jedoch Korrelationen:

\begin{itemize}[itemsep=2pt]
  \item \textbf{{\"O}lpreis und OPEX}: Bei hohen \"Olpreisen steigen typischerweise auch die Betriebskosten.
  \item \textbf{CAPEX und F\"ordervolumen}: Gr\"o{\ss}ere Projekte erm\"oglichen oft h\"ohere F\"orderung.
  \item \textbf{{\"O}lpreis und CAPEX}: Hochpreisphasen treiben die Infrastrukturkosten nach oben.
\end{itemize}

\textbf{Erweiterungsm\"oglichkeit}: Kopula-basierte Abh\"angigkeitsstrukturen k\"onnen Korrelationen zwischen den Inputvariablen modellieren, ohne die marginalen Verteilungen zu ver\"andern.

\subsection{Fehlende Ex-post-Einfl\"usse und Realoptionen}

Wie in Abschnitt~2.2 dargelegt, modellieren wir eine \glqq Fire-and-Forget\grqq{}-Investition. In der Realit\"at haben Manager die M\"oglichkeit:

\begin{itemize}[itemsep=2pt]
  \item \textbf{Abbruchoption}: Das Projekt bei drohendem Verlust vorzeitig zu stoppen.
  \item \textbf{Erweiterungsoption}: Bei \"uberaus g\"unstiger Entwicklung zus\"atzlich zu investieren.
  \item \textbf{Verz\"ogerungsoption}: Den Investitionszeitpunkt zu verschieben.
\end{itemize}

Diese Realoptionen haben einen positiven Wert, der in unserem Modell nicht ber\"ucksichtigt wird \cite{brennan1985}.

\subsection{Statische Parameter}

Das Modell verwendet feste Verteilungsparameter. In der Praxis k\"onnten sich diese Parameter \"uber die Projektdauer \"andern -- insbesondere der \"Olpreis zeigt empirisch ein \textit{Mean-Reverting}-Verhalten. Al-Harthy \cite{alharthy2007} vergleicht verschiedene stochastische \"Olpreismodelle und zeigt, dass die Modellwahl signifikanten Einfluss auf die Investitionsbewertung hat.

\subsection{Single-Project-Perspektive}

Die Simulation betrachtet ein einzelnes Projekt isoliert. In einem Portfolio-Kontext k\"onnen Diversifikationseffekte auftreten.

% ==================================================================
% ABSCHNITT 6: FAZIT
% ==================================================================
\section{Fazit}

\subsection{Zusammenfassung}

Die Monte-Carlo-Simulation bietet einen signifikanten Mehrwert gegen\"uber deterministischen Investitionsrechnungen in der \"Olindustrie:

\begin{enumerate}[label=\arabic*., itemsep=4pt]
  \item \textbf{Von Einzelpunkt zu Verteilung}: Statt einer einzelnen ROI-Zahl erhalten Entscheidungstr\"ager eine vollst\"andige Risikoverteilung.
  \item \textbf{Wissenschaftlich fundierte Verteilungswahl}: Die Kombination aus Dreiecksverteilungen (f\"ur Expertensch\"atzungen) und Lognormalverteilung (f\"ur Rohstoffpreise) ist datengerecht und theoriekonform.
  \item \textbf{Praktische Entscheidungshilfe}: Value at Risk und Probability of Loss sind direkt entscheidungsrelevante Kennzahlen.
  \item \textbf{Identifikation von Risikotreibern}: Die Sensitivit\"atsanalyse zeigt die relative Bedeutung der Inputvariablen.
\end{enumerate}

\subsection{Ausblick}

Folgende Erweiterungen werden f\"ur die n\"achste Iteration empfohlen:

\begin{itemize}[itemsep=2pt]
  \item \textbf{Korrelationsmodellierung} mittels Copulas zur Erfassung von Abh\"angigkeiten
  \item \textbf{Realoptionsansatz} zur Bewertung von Management-Flexibilit\"at
  \item \textbf{Zeitabh\"angige Parameter} f\"ur dynamische \"Olpreismodelle (Mean-Reversion)
  \item \textbf{Portfolio-Simulation} zur Quantifizierung von Diversifikationseffekten
  \item \textbf{Bayesian Updating} zur Integration neuer Informationen im Projektverlauf
\end{itemize}

% ==================================================================
% LITERATUR
% ==================================================================
\begin{thebibliography}{99}

\bibitem{newendorp2000}
Newendorp, P.\,D.\ \& Schuyler, J.\,R.\ (2000).
\textit{Decision Analysis for Petroleum Exploration}.
2.~Auflage. Planning Press.

\bibitem{brennan1985}
Brennan, M.\,J.\ \& Schwartz, E.\,S.\ (1985).
Evaluating Natural Resource Investments.
\textit{The Journal of Business}, 58(2), 135--157.

\bibitem{alharthy2007}
Al-Harthy, M.\,H.\ (2007).
Stochastic Oil Price Models: Comparison and Impact.
\textit{The Engineering Economist}, 52(3), 269--284.
DOI:~10.1080/00137910701633512.

\bibitem{spe35990}
SPE~35990 (1996).
Probabilistic Drilling-Cost Estimating.
\textit{Society of Petroleum Engineers}.

\bibitem{spe52949}
SPE~52949 (1999).
Comparing Three Methods for Evaluating Oil Projects: Option Pricing, Decision Trees, and Monte Carlo Simulations.
\textit{Society of Petroleum Engineers}.

\bibitem{hertz1964}
Hertz, D.\,B.\ (1964).
Risk Analysis in Capital Investment.
\textit{Harvard Business Review}, 42(1), 95--106.

\bibitem{leonteq2023}
Leonteq Securities AG.\ (2023).
\textit{Leonteq Brent Crude Oil Futures Total Return Index (LEONCOTT Index)}.
Online verf\"ugbar unter: \url{https://indices.leonteq.com/index-details/LEONCOTT\%20Index} (abgerufen am 29.04.2026).

\end{thebibliography}

\vspace{1cm}
\end{document}
